\documentclass[11pt,a4paper]{article}

\usepackage{fontspec}
\usepackage[czech]{babel}
\usepackage{geometry}
\usepackage{xcolor}
\usepackage{colortbl}
\usepackage{tabularx}
\usepackage{enumitem}
\usepackage{titlesec}
\usepackage{hyperref}
\usepackage{array}
\usepackage{microtype}

\geometry{
  left=22mm,
  right=22mm,
  top=20mm,
  bottom=22mm
}

\setmainfont{Arial}
\setsansfont{Arial}
\hypersetup{
  colorlinks=true,
  linkcolor=black,
  urlcolor=blue
}

\definecolor{AuraDark}{HTML}{1F2937}
\definecolor{AuraMuted}{HTML}{4B5563}
\definecolor{AuraLine}{HTML}{D1D5DB}
\definecolor{AuraSoft}{HTML}{F3F4F6}

\titleformat{\section}
  {\Large\bfseries\color{AuraDark}}
  {}
  {0pt}
  {}
  [\vspace{2pt}{\color{AuraLine}\titlerule}]

\titleformat{\subsection}
  {\normalsize\bfseries\color{AuraDark}}
  {}
  {0pt}
  {}

\setlist[itemize]{
  topsep=3pt,
  itemsep=2pt,
  parsep=0pt,
  leftmargin=18pt
}

\setlength{\parindent}{0pt}
\setlength{\parskip}{6pt}
\renewcommand{\arraystretch}{1.25}

\newcommand{\labelrow}[2]{\textbf{#1} & #2 \\}

\begin{document}
\pagecolor{white}
\color{black}

\begin{center}
  {\LARGE \textbf{Progress SP}}\\[3pt]
  {\large Aktuální stav projektu AURA}\\[8pt]
  {\color{AuraMuted} David Strnadel, AK8BO - Bezpečnost operačních systémů}\\
  {\color{AuraMuted} Aktualizace: 12. 5. 2026 \quad | \quad obhajoba: pátek 15. 5. 2026}
\end{center}

\vspace{4pt}

\begin{tabularx}{\textwidth}{>{\bfseries}p{36mm}X}
Hlavní výstup & Strukturované bezpečnostní a QA reporty nad daty dostupnými běžné Android aplikaci \\
Stav repozitáře & Aktuální implementace je přiložená v GitHub repozitáři projektu \\
\end{tabularx}

\section{Co je cílem projektu}

Na začátku projektu jsem uvažoval hlavně o mobilním bezpečnostním agentovi pro Android. Během vývoje se ale ukázalo, že nejzajímavější a zároveň realističtější směr není vytvořit další aplikaci pro běžného uživatele ani konkurovat antivirovým řešením nebo Google Play Protect.

Aktuální cíl projektu je jiný: AURA má být výzkumný a B2B reporting nástroj. Android aplikace v tomto pojetí není finální produkt sama o sobě. Slouží hlavně jako sběrná sonda, která z no-root prostředí získá maximum relevantních dat, uloží je do strukturovaného exportu a nad nimi se následně generují reporty.

Produktem projektu jsou tedy reporty, ne uživatelské UI aplikace. Cílovým uživatelem nejsou běžní uživatelé telefonu, ale spíše firmy a týmy zabývající se Android vývojem, QA, bezpečnostní revizí aplikací nebo interním hodnocením rizik. Smyslem je dodat jim jiný pohled na aplikaci: ne jako pentest, ne jako malware analýzu, ale jako vysvětlitelný report založený na datech, ke kterým se dá dostat bez root oprávnění.

Praktická otázka projektu tedy zní:

\begin{quote}
\textit{Kam až se dá v Androidu dostat bez root oprávnění a co užitečného lze z takto získaných dat říci o aplikacích, jejich oprávněních, komponentách, původu, rizikových kombinacích a obranné ploše?}
\end{quote}

\section{Aktuální stav implementace}

Projekt je nyní ve fázi funkčního research MVP. V repozitáři je implementovaná Android aplikace, která sbírá dostupná data o aplikacích v zařízení a vytváří strukturovaný JSON export. Nad exportem pak existují podpůrné nástroje pro evaluaci, redakci citlivých údajů a generování reportů.

Android část umí sbírat zejména:

\begin{itemize}
  \item základní metadata aplikací dostupných přes \texttt{PackageManager},
  \item deklarovaná a udělená oprávnění,
  \item aktivity, služby, receivery a providery,
  \item exportované komponenty a informaci, zda jsou chráněné oprávněním,
  \item instalační zdroj, podpisové digesty a systémový původ aplikace,
  \item vybrané special-access signály, například Accessibility Service, Notification Listener, overlay a request-install-packages,
  \item omezenou UsageStats korelaci v research scénáři, pokud je explicitně povolená.
\end{itemize}

Nad těmito daty běží rozhodovací a vysvětlovací vrstva. Ta se nesnaží aplikace hodnotit jen podle počtu oprávnění. Rozlišuje například capability exposure, legitimitu role aplikace, provenienci, konkrétní evidence, akceschopnost uživatele a nejistotu.

Výsledek se ukládá jako JSON export a je možné z něj generovat reporty. Právě reporty jsou pro aktuální směr projektu hlavní výstup.

\section{Dosavadní výstupy}

V projektu už mám připravené tyto části:

\begin{itemize}
  \item Android research aplikaci AURA v Kotlinu.
  \item Datový model pro snapshoty aplikací, evidence items, risk vector, decision trace, user risk story a evidence graph.
  \item Rozhodovací logiku pro role inference, provenance classification, risk decision engine, temporal episodes a defensive-surface findings.
  \item JSON export výsledků do privátního úložiště aplikace.
  \item Sadu neškodných testovacích aplikací pro řízené scénáře.
  \item Python evaluator pro porovnání AURA výstupů proti jednodušším baseline modelům.
  \item Report generator pro Markdown/HTML reporty.
  \item Export redactor pro přípravu reportů v různých úrovních citlivosti.
  \item Dokumentaci k observability limitům, risk modelu, privacy/ethics a testovacím scénářům.
\end{itemize}

Lokálně jsem ověřil dostupné testy, které nevyžadují běžící Android emulator. Gradle testovací běh pro \texttt{researchFullStandardDebug} unit testy skončil úspěšně. Dále jsem spustil Python unit testy pro podpůrné nástroje, kde proběhlo 27 testů a všechny skončily úspěšně.

\section{Reporty jako hlavní produkt}

Nejdůležitější posun oproti původní představě je v tom, jak chápu výstup projektu. UI Android aplikace je zatím velmi jednoduché a neplánuji ho do obhajoby výrazně zlepšovat. Není to priorita, protože pro aktuální použití nemá být hlavní hodnotou klikací aplikace.

Hlavní hodnotou jsou reporty, které mohou být užitečné například pro:

\begin{itemize}
  \item QA tým, který chce další pohled na Android aplikaci před vydáním,
  \item vývojářský tým, který chce rychle vidět podezřelé kombinace oprávnění a komponent,
  \item bezpečnostní review, které potřebuje oddělit konkrétní zjištění od pouhé nejistoty,
  \item interní audit aplikací, kde je důležité pracovat s omezeními no-root prostředí.
\end{itemize}

Report nemá říkat „aplikace je malware“. Má spíš říct: toto bylo pozorováno, toto pozorováno nebylo, tyto části dávají smysl vzhledem k roli aplikace, tyto části si zaslouží revizi a tady jsou limity toho, co no-root agent vůbec může zjistit.

\section{Co zbývá do obhajoby}

Do obhajoby už nechci projekt rozšiřovat o nové velké funkce. Priorita je uzavřít existující MVP, spustit finální scénáře, vybrat ukázkové výstupy a připravit prezentaci tak, aby byl projekt pochopitelný i bez detailní znalosti celého repozitáře.

\begin{tabularx}{\textwidth}{>{\bfseries}p{34mm}X}
\rowcolor{AuraSoft}
Datum & Plán dokončení \\
Úterý 12. 5. 2026 & Odevzdat progress SP, sjednotit popis projektu a uzavřít rozsah pro obhajobu. \\
Středa 13. 5. 2026 & Spustit finální emulátorové scénáře, stáhnout JSON exporty a uložit výsledky evaluatoru. \\
Čtvrtek 14. 5. 2026 & Zpracovat výsledky do stručné tabulky, vybrat ukázkové reporty a dokončit prezentaci. \\
Pátek 15. 5. 2026 & Projít demo, připravit repozitář a obhájit projekt. \\
\end{tabularx}

\section{Další směr po odevzdání}

Ve vývoji chci pokračovat i po obhajobě. No-root část beru jako první výzkumnou linii: ukazuje, co lze zjistit z pozice běžné aplikace. Vedle toho bych chtěl časem přidat samostatný root lab nebo laboratorní režim, který by umožnil sbírat hlubší systémová data v kontrolovaném prostředí.

Tím by šlo porovnat dvě úrovně pohledu:

\begin{itemize}
  \item co je možné říci z čistě no-root dat,
  \item a co se změní, když má výzkumný lab hlubší přístup k systému.
\end{itemize}

Tento směr by mohl být užitečný pro další výzkum i pro B2B reporting. Firmám by bylo možné ukázat nejen samotný nález, ale také rozdíl mezi běžně dostupnou observabilitou a hlubší laboratorní analýzou.

\section{Zhodnocení}

Aktuální stav projektu odpovídá závěrečné fázi původního plánu. Rešeršní a návrhová část se postupně přetavila do funkčního MVP, které sbírá dostupná Android data, vytváří exporty a umožňuje z nich generovat reporty.

Za hlavní přínos projektu považuji právě tento posun: nesnažím se dodat další mobilní bezpečnostní aplikaci pro koncového uživatele, ale ověřuji, jak daleko se dá dostat s no-root sběrem dat a jak z těchto dat vytvořit srozumitelný, opatrný a prakticky použitelný report pro Android vývoj a QA.

\end{document}
